\documentclass[11pt]{article}
\usepackage{fontspec}
\setmainfont[
  Path=fonts/,
  Extension=.ttf,
  UprightFont=*-400,
  BoldFont=*-700
]{NotoSerifSC}
\setsansfont[
  Path=fonts/,
  Extension=.ttf,
  UprightFont=*-400,
  BoldFont=*-700
]{NotoSerifSC}
\XeTeXlinebreaklocale "zh"
\XeTeXlinebreakskip=0pt plus 1pt
\usepackage[a4paper,margin=1.70cm]{geometry}
\usepackage{amsmath,amssymb,mathtools,bm}
\usepackage{booktabs,longtable,array}
\usepackage{enumitem}
\setlength{\parindent}{2em}
\setlength{\parskip}{0.35em}
\setlist{nosep,leftmargin=2em}
\allowdisplaybreaks[3]
\numberwithin{equation}{section}
\pagestyle{plain}

\newcommand{\R}{\mathbb{R}}
\newcommand{\E}{\mathbb{E}}
\newcommand{\Pp}{\mathbb{P}}
\newcommand{\1}{\mathbf{1}}
\newcommand{\T}{\mathsf{T}}
\newcommand{\F}{\mathrm{F}}
\newcommand{\cN}{\mathcal{N}}
\newcommand{\cL}{\mathcal{L}}
\newcommand{\cS}{\mathcal{S}}
\newcommand{\cH}{\mathcal{H}}
\newcommand{\cA}{\mathcal{A}}
\newcommand{\cM}{\mathcal{M}}
\newcommand{\cT}{\mathcal{T}}
\newcommand{\diag}{\operatorname{diag}}
\newcommand{\tr}{\operatorname{tr}}
\newcommand{\Var}{\operatorname{Var}}
\newcommand{\Cov}{\operatorname{Cov}}
\newcommand{\softmax}{\operatorname{softmax}}
\newcommand{\KL}{D_{\mathrm{KL}}}
\newcommand{\sg}{\operatorname{sg}}
\begin{document}
    \begin{center}
      {\LARGE\bfseries 34\_Toy Models of Superposition}
    \end{center}
    \vspace{0.8em}

    \section{符号与维度}
    \begin{longtable}{>{$}l<{$} >{$}l<{$} p{10cm}}
    \toprule
    \text{符号} & \text{维度} & \text{含义}\\
    \midrule
    x,\hat x & \R^n & 稀疏输入特征与重建。\\
        W & \R^{m\times n} & $m<n$ 的线性瓶颈。\\
        w_i & \R^m & 第 $i$ 个特征的表示方向。\\
        I_i & \R_{>0} & 特征重要度。\\
        G=W^{\T} W & \R^{n\times n} & 表示方向 Gram 矩阵。\\
    \bottomrule
    \end{longtable}

    所有向量默认是列向量；未特别说明时，期望同时对数据、噪声和采样时刻取值。下文只展开论文中承担方法逻辑的公式，并把所需假设写在推导发生的位置。

\section{线性自编码器中的干扰展开}

设 $n$ 个理想特征 $x\in\R^n$ 被压入 $m<n$ 维，列向量
$w_i\in\R^m$ 组成 $W=[w_1,ldots,w_n]$。绑权重线性重建
\begin{equation}
 \widehat x=W^TWx.
\end{equation}
第 $i$ 个坐标
\begin{equation}
 \widehat x_i=w_i^T\sum_jw_jx_j
 =\|w_i\|_2^2x_i+\sum_{j\ne i}(w_i^Tw_j)x_j.
\end{equation}
第一项是自信号，第二项是其他特征通过 Gram 非对角元产生的干扰。

令重要性权重 $I_i\ge0$，平方损失
\begin{equation}
 \cL=\E\sum_iI_i(x_i-\widehat x_i)^2.
\end{equation}
假设各特征独立、均值零、方差 $v_j$，交叉项期望消失，得到
\begin{align}
 \E(x_i-\widehat x_i)^2
 &=(1-\|w_i\|^2)^2v_i
 +\sum_{j\ne i}(w_i^Tw_j)^2v_j.
\end{align}
所以
\begin{equation}
 \boxed{\cL=\sum_iI_i(1-\|w_i\|^2)^2v_i
 +\sum_i\sum_{j\ne i}I_i(w_i^Tw_j)^2v_j.}
\end{equation}
正交表示消除干扰但最多容纳 $m$ 个非零方向；超位置用小而非零的内积容纳更多特征。

\section{稀疏性为何降低期望干扰}

若
\begin{equation}
 x_j=\begin{cases}
 0,&\text{概率 }S_j,\\
 \xi_j,&\text{概率 }1-S_j,
 \end{cases}
\end{equation}
且 $\E\xi_j=0$、$\E\xi_j^2=\nu_j$，则
\begin{equation}
 v_j=\E x_j^2=(1-S_j)\nu_j.
\end{equation}
干扰项变为
\begin{equation}
 \sum_{j\ne i}(w_i^Tw_j)^2(1-S_j)\nu_j.
\end{equation}
当特征很少同时激活，方向非正交造成的平均代价按激活概率下降；
这给“稀疏特征可以超位置”一个定量解释。

若特征均值不为零，$\E[x_jx_k]$ 的交叉项不再消失，损失还包含
\begin{equation}
 \sum_{j\ne k}(w_i^Tw_j)(w_i^Tw_k)\E[x_jx_k].
\end{equation}
因此独立零均值假设是上述简化式的重要前提。

\section{Gram 矩阵与 Welch 下界}

假设 $n$ 个方向单位化，$\|w_i\|=1$。Gram 矩阵
$G=W^TW\in\R^{n\times n}$，秩至多 $m$，迹
\begin{equation}
 \tr G=\sum_i\|w_i\|^2=n.
\end{equation}
设 $G$ 的非零特征值为 $\lambda_1,ldots,\lambda_m$。由 Cauchy--Schwarz，
\begin{equation}
 \|G\|_F^2=\sum_{k=1}^{m}\lambda_k^2
 \ge\frac1m\left(\sum_k\lambda_k\right)^2
 =\frac{n^2}{m}.
\end{equation}
另一方面
\begin{equation}
 \|G\|_F^2=\sum_iG_{ii}^2+\sum_{i\ne j}G_{ij}^2
 =n+\sum_{i\ne j}(w_i^Tw_j)^2.
\end{equation}
故
\begin{equation}
 \boxed{\frac1{n(n-1)}\sum_{i\ne j}(w_i^Tw_j)^2
 \ge\frac{n-m}{m(n-1)}.}
\end{equation}
于是最大相干性也满足
\begin{equation}
 \max_{i\ne j}|w_i^Tw_j|
 \ge\sqrt{\frac{n-m}{m(n-1)}}.
\end{equation}
这说明 $n>m$ 时不可能让全部方向严格正交，最优只能均匀分配干扰。

\section{ReLU 输出与负干扰的非对称性}

带偏置的模型
\begin{equation}
 \widehat x=\operatorname{ReLU}(W^TWx+b)
\end{equation}
在第 $i$ 坐标预激活
\begin{equation}
 a_i=\|w_i\|^2x_i+\sum_{j\ne i}(w_i^Tw_j)x_j+b_i.
\end{equation}
梯度门为
\begin{equation}
 \frac{\partial\widehat x_i}{\partial a_i}
 =\1\{a_i>0\}
\end{equation}
（零点取任一子梯度）。如果真实特征非负，适当负偏置可让小的纯干扰预激活被截为零，
而强自信号仍越过阈值。这使正负内积的代价不再对称，几何最优可形成多边形、反向对等相位。

\section{单特征表示、超位置与忽略的比较}

考虑一个重要性 $I$、激活概率 $p=1-S$ 的特征。
若完全忽略，输出恒为零，代价
\begin{equation}
 L_{\rm ignore}=I\E x^2=Ip\nu.
\end{equation}
若独占单位方向且解码无偏，理想代价接近零。若加入超位置，它获得自信号但产生/承受干扰，
粗略代价可写
\begin{equation}
 L_{\rm super}\approx I(1-\|w_i\|^2)^2p\nu
 +Ip\sum_{j\ne i}p_j\nu_j(w_i^Tw_j)^2.
\end{equation}
随着 $I$ 或 $p$ 增加，忽略代价线性上升，模型更愿分配专属维度；
当 $p$ 很小，干扰又被联合激活概率削弱，超位置更划算。不同方案损失曲线交叉时，
全局最优结构会突然切换，形成观察到的相变。

\section{多义神经元与可解释性的区别}

隐藏坐标 $h=Wx$ 的第 $k$ 个神经元
\begin{equation}
 h_k=\sum_iW_{ki}x_i
\end{equation}
通常同时响应多个特征；但表示方向是列 $w_i$，不必与坐标轴对齐。
若 $R\in\R^{m\times m}$ 正交，替换
\begin{equation}
 W\leftarrow RW
\end{equation}
不会改变 Gram 矩阵
\begin{equation}
 (RW)^T(RW)=W^TW,
\end{equation}
因而重建函数不变，却会完全旋转神经元坐标。可辨识的是特征方向间的几何关系，
不是某个隐藏坐标的绝对语义。

\section{压缩感知类比的边界}

若 $x$ 是 $k$-稀疏，随机测量 $h=Wx$ 在满足受限等距性质时可由非线性优化恢复。
RIP 要求所有 $k$-稀疏向量满足
\begin{equation}
 (1-\delta)\|x\|_2^2\le\|Wx\|_2^2\le(1+\delta)\|x\|_2^2.
\end{equation}
随机矩阵通常需 $m=O(k\log(n/k)/\delta^2)$。玩具模型解码器
$\operatorname{ReLU}(W^T h+b)$ 并不是通用 $\ell_1$ 恢复算法，因此“超位置类似压缩感知”是几何动机，
不能直接继承精确恢复定理。

\section{低秩更新的完整矩阵微分}

令冻结权重 $W_0\in\R^{m\times n}$，低秩更新
\begin{equation}
 W=W_0+cBA,\qquad
 B\in\R^{m\times r},\quad A\in\R^{r\times n},\quad c=\alpha/r.
\end{equation}
对一个输入 $x\in\R^n$，
\begin{equation}
 h=W_0x+cB(Ax).
\end{equation}
令 $y=Ax\in\R^r$、输出梯度 $g=\partial\cL/\partial h\in\R^m$。
微分为
\begin{equation}
 dh=c\,dB\,y+cB\,dA\,x+cBA\,dx+W_0dx.
\end{equation}
由 $d\cL=g^Tdh$ 以及
$a^TdMb=\langle ab^T,dM\rangle_\F$，逐项得到
\begin{align}
 \boxed{\nabla_B\cL=c,g,y^T=c,g,x^TA^T,}\\
 \boxed{\nabla_A\cL=c,B^Tg,x^T,}\\
 \boxed{\nabla_x\cL=W_0^Tg+cA^TB^Tg.}
\end{align}
若批输入按行堆成 $X\in\R^{b\times n}$，上游梯度为
$G\in\R^{b\times m}$，则
\begin{align}
 \nabla_B\cL&=cG^TXA^T,\\
 \nabla_A\cL&=cB^TG^TX.
\end{align}

\subsection{rank-1 专家分解}

写 $B=[b_1,\ldots,b_r]$、$A$ 的第 $i$ 行为 $a_i^T$，则
\begin{equation}
 BA=\sum_{i=1}^{r}b_i a_i^T,
 \qquad BAx=\sum_{i=1}^{r}b_i(a_i^Tx).
\end{equation}
每一项均为秩至多 $1$ 的更新。对单样本，
\begin{equation}
 \nabla_{b_i}\cL=c(a_i^Tx)g,\qquad
 \nabla_{a_i}\cL=c(b_i^Tg)x.
\end{equation}
于是列 $b_i$ 的更新强度由标量路由量 $a_i^Tx$ 控制；这也是把 rank 分量解释为隐式专家的数学依据。

\section{初始化时两因子的非对称动力学}

常见初始化为 $B_0=0$、$A_0$ 随机。此时
\begin{equation}
 \Delta W_0=cB_0A_0=0,
\end{equation}
所以前向不改变基座。梯度却满足
\begin{equation}
 \nabla_B\cL\big|_{0}=cG^TXA_0^T,\qquad
 \nabla_A\cL\big|_{0}=cB_0^TG^TX=0.
\end{equation}
第一步梯度下降
\begin{equation}
 B_1=-\eta cG^TXA_0^T,\qquad A_1=A_0,
\end{equation}
因此第二步开始 $A$ 才获得非零梯度。若反过来令 $A_0=0$、$B_0$ 随机，
则第一步只有 $A$ 更新。两种初始化前向等价，但早期优化动力学不同。

还有尺度不识别性：任意 $s\ne0$ 满足
\begin{equation}
 BA=(sB)(s^{-1}A).
\end{equation}
然而带独立权重衰减时
\begin{equation}
 \|sB\|_F^2+\|s^{-1}A\|_F^2
 =s^2\|B\|_F^2+s^{-2}\|A\|_F^2,
\end{equation}
最优平衡尺度由求导得到
\begin{equation}
 s^4=\frac{\|A\|_F^2}{\|B\|_F^2}.
\end{equation}
故因子正则化会隐式选择分解，即使乘积 $BA$ 相同。

\section{最佳秩约束近似}

设目标全量更新的奇异值分解为
\begin{equation}
 \Delta W^*=U\Sigma V^T
 =\sum_{i=1}^{q}\sigma_i u_iv_i^T,
 \qquad\sigma_1\ge\cdots\ge\sigma_q\ge0.
\end{equation}
截断矩阵
\begin{equation}
 \Delta W_r=\sum_{i=1}^{r}\sigma_i u_iv_i^T
\end{equation}
满足 Eckart--Young--Mirsky 结论
\begin{equation}
 \min_{\operatorname{rank}(M)\le r}
 \|\Delta W^*-M\|_F^2
 =\sum_{i>r}\sigma_i^2.
\end{equation}
一种直接证明如下。因 Frobenius 范数在正交变换下不变，令
$C=U^TMV$，则 $\operatorname{rank}(C)\le r$ 且
\begin{equation}
 \|\Delta W^*-M\|_F^2=\|\Sigma-C\|_F^2.
\end{equation}
任意秩不超过 $r$ 的 $C$ 至多精确保留 $r$ 个相互正交奇异方向；
由奇异值的变分刻画或 Pythagoras 分解，剩余平方能量至少为
$\sum_{i>r}\sigma_i^2$。取
$C=\diag(\sigma_1,\ldots,\sigma_r,0,\ldots)$ 达到该下界。

谱范数版本为
\begin{equation}
 \min_{\operatorname{rank}(M)\le r}\|\Delta W^*-M\|_2=\sigma_{r+1}.
\end{equation}
这只是表达能力上界；训练能否找到截断 SVD 解还取决于非凸因子化、数据和优化器。

\section{子空间限制与梯度投影}

若 $A$ 冻结，只训练 $B$，可达更新集合为
\begin{equation}
 \cS_A=\{BA:B\in\R^{m\times r}\}.
\end{equation}
其中每一行都位于 $A$ 的行空间。若 $A$ 满行秩，其行空间正交投影为
\begin{equation}
 P_A=A^T(AA^T)^{-1}A.
\end{equation}
给定任意目标更新 $M$，最小二乘问题
\begin{equation}
 \min_B\|M-BA\|_F^2
\end{equation}
的正规方程为
\begin{equation}
 (BA-M)A^T=0,\qquad
 B^*=MA^T(AA^T)^{-1}.
\end{equation}
故最佳可达更新
\begin{equation}
 B^*A=MP_A,\qquad
 \min_B\|M-BA\|_F^2=\|M(I-P_A)\|_F^2.
\end{equation}
这把“冻结随机下投影”的表达损失精确写成目标在随机行空间外的能量。

\section{随机投影的尺度与集中}

设 $A_{ij}$ 独立、均值零、方差 $1/r$。对固定 $x\in\R^n$，
\begin{align}
 \E\|Ax\|_2^2
 &=\sum_{i=1}^{r}\E\left(\sum_jA_{ij}x_j\right)^2\\
 &=\sum_i\sum_j\E[A_{ij}^2]x_j^2
 =\|x\|_2^2.
\end{align}
若 $A_{ij}\sim\cN(0,1/r)$，则
\begin{equation}
 \frac{\|Ax\|_2^2}{\|x\|_2^2}
 \sim\frac1r\chi_r^2.
\end{equation}
利用卡方尾界，对 $0<\varepsilon<1$，
\begin{equation}
 \Pp\left(
 \left|\|Ax\|_2^2-\|x\|_2^2\right|
 \ge\varepsilon\|x\|_2^2\right)
 \le2\exp\left(-\frac r8\varepsilon^2\right).
\end{equation}
对有限集合 $N$ 个向量做并合界，若
\begin{equation}
 r\ge\frac8{\varepsilon^2}\log\frac{2N}{\delta},
\end{equation}
则所有向量同时保持范数的概率至少为 $1-\delta$。对点对差向量应用同一结论即得到距离保持。

稀疏随机投影若每个元素以概率 $\rho$ 非零，则为保持同样二阶矩，非零幅值应按
$1/\sqrt{\rho r}$ 缩放：
\begin{equation}
 A_{ij}=\begin{cases}
 +1/\sqrt{\rho r},&\rho/2,\\
 -1/\sqrt{\rho r},&\rho/2,\\
 0,&1-\rho.
 \end{cases}
\end{equation}
此时 $\E[A_{ij}^2]=1/r$，但四阶矩为
$\E[A_{ij}^4]=1/(\rho r^2)$；$\rho$ 越小，尾部越重，获得同样集中精度通常需要更大 $r$。

\section{任务更新的干扰量}

两个任务的更新 $U_i=B_iA_i$、$U_j=B_jA_j$ 的 Frobenius 内积为
\begin{align}
 \langle U_i,U_j\rangle_F
 &=\tr((B_iA_i)^TB_jA_j)\\
 &=\tr(B_i^TB_jA_jA_i^T).
\end{align}
因此
\begin{equation}
 |\langle U_i,U_j\rangle_F|
 \le\|B_i^TB_j\|_F\|A_jA_i^T\|_F.
\end{equation}
随机下投影近似正交只控制第二个因子；若 $B_i,B_j$ 范数无界，绝对干扰仍可任意大。
对加权合并 $U=\sum_iw_iU_i$，
\begin{equation}
 \|U\|_F^2
 =\sum_iw_i^2\|U_i\|_F^2
 +2\sum_{i<j}w_iw_j\langle U_i,U_j\rangle_F.
\end{equation}
交叉项就是任务干扰对合并范数的精确贡献，而不是抽象的“方向冲突”。

\section{参数量、训练内存与合并}

全量矩阵参数量为 $mn$，低秩因子参数量为
\begin{equation}
 r(m+n),
\end{equation}
压缩要求 $r<mn/(m+n)$。训练时只有可训练因子需要梯度和优化器状态；
冻结基座仍需存储权重并参与前向、反向的输入梯度计算。

推理前可合并
\begin{equation}
 W_{\rm merged}=W_0+cBA.
\end{equation}
合并后对单个适配器不增加矩阵乘次数，但丢失了动态切换不同适配器的便利。
若基座以低比特存储，通常必须先在计算精度中反量化再加 $BA$；直接在整数码空间相加并不等价。

\section{条件期望作为 $L^2$ 正交投影}

令随机向量 $Y\in L^2$，观测 $Z$ 生成子空间
\begin{equation}
 \mathcal H_Z=\{f(Z):\E\|f(Z)\|^2<\infty\}.
\end{equation}
条件期望 $m(Z)=\E[Y\mid Z]$ 满足任意 $h(Z)\in\mathcal H_Z$：
\begin{equation}
 \E[(Y-m(Z))^Th(Z)]=0.
\end{equation}
因为
\begin{align}
 \E[(Y-m)^Th]
 &=\E\{\E[(Y-m)^Th\mid Z]\}\\
 &=\E\{h^T\E[Y-m\mid Z]\}=0.
\end{align}
于是 Pythagoras 恒等式
\begin{equation}
 \E\|Y-f(Z)\|^2
 =\E\|Y-m(Z)\|^2+\E\|m(Z)-f(Z)\|^2.
\end{equation}
多数回归、去噪、流匹配和局部目标的“最优预测器”都由这一投影结论得到。

全方差公式也是同一正交分解：
\begin{equation}
 \Var(Y)=\E\Var(Y\mid Z)+\Var(\E[Y\mid Z]).
\end{equation}
第一项是观测 $Z$ 后仍不可约的不确定性，第二项是不同 $Z$ 条件均值的变化。

\section{Jensen、对数和与变分界}

对凹函数 $\log$，
\begin{equation}
 \log\E X\ge\E\log X,
 \qquad X>0.
\end{equation}
差距可以写成 KL。令未归一权重 $w(z)>0$、提议分布 $q(z)$，
配分函数 $Z=\E_qw(z)$，定义
\begin{equation}
 p^*(z)=\frac{q(z)w(z)}{Z}.
\end{equation}
则
\begin{align}
 \log Z-\E_q\log w
 &=\E_q\log\frac{Z}{w(z)}\\
 &=\E_q\log\frac{q(z)}{p^*(z)}\\
 &=D_{\rm KL}(q\|p^*)\ge0.
\end{align}
因此 Jensen 下界何时紧不只取决于“方差小”，而是提议分布是否等于归一化后的目标分布。

log-sum-exp
\begin{equation}
 \operatorname{LSE}(x)=\log\sum_{i=1}^{n}e^{x_i}
\end{equation}
满足
\begin{equation}
 \max_i x_i\le\operatorname{LSE}(x)
 \le\max_i x_i+\log n.
\end{equation}
减去最大值 $m=\max_ix_i$ 后计算
\begin{equation}
 \operatorname{LSE}(x)=m+\log\sum_ie^{x_i-m}
\end{equation}
可避免指数上溢且数学等价。

\section{Hoeffding 不等式的矩母函数推导}

设独立 $X_i\in[a_i,b_i]$，$S=\sum_i(X_i-\E X_i)$。
对任意 $\lambda>0$，Markov 不等式
\begin{equation}
 \Pp(S\ge t)=\Pp(e^{\lambda S}\ge e^{\lambda t})
 \le e^{-\lambda t}\E e^{\lambda S}.
\end{equation}
独立性给出矩母函数乘积。Hoeffding 引理
\begin{equation}
 \E e^{\lambda(X_i-\E X_i)}
 \le\exp\left(\frac{\lambda^2(b_i-a_i)^2}{8}\right).
\end{equation}
故
\begin{equation}
 \Pp(S\ge t)
 \le\exp\left(-\lambda t+\frac{\lambda^2}{8}
 \sum_i(b_i-a_i)^2\right).
\end{equation}
对 $\lambda$ 最小化，
\begin{equation}
 \lambda^*=\frac{4t}{\sum_i(b_i-a_i)^2},
\end{equation}
得到
\begin{equation}
 \Pp(S\ge t)
 \le\exp\left(-\frac{2t^2}{\sum_i(b_i-a_i)^2}\right).
\end{equation}
对下尾应用于 $-S$，再 union bound 得双侧界。

\section{Bernstein 界利用方差信息}

若独立零均值 $X_i$ 满足 $|X_i|\le M$，总方差
$v=\sum_i\E X_i^2$，Bernstein 界
\begin{equation}
 \Pp\left(\sum_iX_i\ge t\right)
 \le\exp\left(-\frac{t^2}{2(v+Mt/3)}\right).
\end{equation}
当 $t\ll v/M$ 时近似高斯尾 $e^{-t^2/(2v)}$；
当 $t$ 很大时转为 $e^{-3t/(2M)}$。相比只使用取值区间的 Hoeffding，
方差远小于最坏范围时 Bernstein 更紧。

对 Bernoulli 均值 $\widehat p=N^{-1}\sum_iX_i$，$v=Np(1-p)$，
\begin{equation}
 \Pp(|\widehat p-p|\ge\epsilon)
 \le2\exp\left(-\frac{N\epsilon^2}
 {2[p(1-p)+\epsilon/3]}\right).
\end{equation}

\section{并合界与同时保证}

事件 $E_1,…,E_M$ 不要求独立也有
\begin{equation}
 \Pp\left(\bigcup_{j=1}^{M}E_j\right)
 \le\sum_{j=1}^{M}\Pp(E_j).
\end{equation}
若每个失败概率至多 $\delta/M$，则全部结论同时成立的概率至少 $1-\delta$。
将单点集中界扩展到有限候选集合通常产生 $\log M$ 项。例如
\begin{equation}
 2M e^{-2N\epsilon^2}\le\delta
\end{equation}
等价于
\begin{equation}
 \epsilon\ge\sqrt{\frac{\log(2M/\delta)}{2N}}.
\end{equation}
若候选集合无限，需要覆盖数、VC 维或 Rademacher 复杂度，而不能直接把 $M=\infty$ 代入。

\section{Delta 方法与非线性估计误差}

设 $\sqrt N(\widehat\theta-\theta)\Rightarrow\cN(0,\Sigma)$，
光滑函数 $g$ 的一阶展开
\begin{equation}
 g(\widehat\theta)=g(\theta)+
 \nabla g(\theta)^T(\widehat\theta-\theta)+o_p(N^{-1/2}).
\end{equation}
于是
\begin{equation}
 \sqrt N[g(\widehat\theta)-g(\theta)]
 \Rightarrow\cN(0,\nabla g^T\Sigma\nabla g).
\end{equation}
例如 $g(p)=\log p$，方差近似
\begin{equation}
 \Var(\log\widehat p)\approx\frac{\Var(\widehat p)}{p^2}.
\end{equation}
当 $p$ 很小时误差被 $1/p^2$ 放大，解释了稀有事件对数概率和概率比估计的不稳定。

二阶展开还给非线性偏差
\begin{equation}
 \E g(\widehat\theta)-g(\theta)
 \approx\frac12\tr(H_g(\theta)\Cov(\widehat\theta)).
\end{equation}
FID、比率、倒数和归一化优势等即使基础估计无偏，非线性组合也可有有限样本偏差。

\section{重要性采样与有效样本量}

目标期望 $\E_p f(X)$ 用提议 $q$ 采样，权重
$w(x)=p(x)/q(x)$：
\begin{equation}
 \widehat\mu=\frac1N\sum_{i=1}^{N}w(X_i)f(X_i),
 \qquad X_i\sim q.
\end{equation}
若 $p$ 的归一化常数未知，使用自归一估计
\begin{equation}
 \widetilde\mu=\frac{\sum_iw_if_i}{\sum_iw_i},
\end{equation}
它一般有 $O(1/N)$ 偏差，但常更实用。权重退化可用
\begin{equation}
 N_{\rm eff}=\frac{(\sum_iw_i)^2}{\sum_iw_i^2}
\end{equation}
衡量；权重相等时为 $N$，单个权重主导时接近 $1$。

比率裁剪 $\widetilde w=\min(w,c)$ 降低方差，但引入偏差
\begin{equation}
 \E_q[(w-\widetilde w)f]
 =\E_q[(w-c)_+f].
\end{equation}
若 $|f|\le F$，绝对偏差至多
\begin{equation}
 F\E_q[(w-c)_+].
\end{equation}

\section{校准与 Brier 分解}

二元预测概率 $P\in[0,1]$、标签 $Y\in\{0,1\}$。条件真实频率
$m(P)=\E[Y\mid P]$。Brier 分数
\begin{align}
 \E(P-Y)^2
 &=\E(P-m(P))^2+\E(m(P)-Y)^2,
\end{align}
交叉项因条件期望为零而消失。第一项是校准误差，第二项是给定预测分组后的不可约噪声。
若模型完美校准，$m(P)=P$，第一项为零，但单次预测仍可错误；
校准不是零错误率，而是概率与长期频率一致。

\end{document}
